% ID: SEL-D0032-OA-E04
% ESTADO: cotejada. Segunda pasada visual y borrador compilado revisados.
% Fuente: Selectividad/matematicas_aplicadas/titular_junio_EXAMEN 5-ANDALUCÍA.pdf
% SHA-256: 5f24b200fb6cf34f6d72016ec3dc35787b4074de9b7752677e5b0d959ed0dfcd
% Página: 1; opción A; ejercicio original 4.

La altura de los estudiantes de 2\textsuperscript{o} de bachillerato de un centro sigue una ley Normal de media 165 cm y desviación típica 10 cm.
\begin{enumerate}[label=\alph*)]
\item \textbf{(1 punto)} \textquestiondown{}Qué distribución sigue la altura media de las muestras de tamaño 25?
\item \textbf{(1.5 puntos)} Se elige al azar una muestra de 25 estudiantes y se les mide la altura. \textquestiondown{}Cuál es la probabilidad de que la altura media de esa muestra supere 160 cm?
\end{enumerate}
